\subsection*{一步一步建立直观}

\paragraph{1. 先承认测量需要一段距离}

把 $f(x)=x^2$ 当作位置读数。在时刻 $a$ 和 $a+h$ 之间，平均速度由两点连成的割线斜率给出。拖动第二点，让 $h$ 保持非零；此时分子、分母都来自真实的两次测量，没有任何一步要求除以零。



\paragraph{2. 缩小窗口，看什么保持稳定}

展开平方可得割线斜率 $2a+h$。缩小 $h$ 的同时，位移也在缩小，但它们的比值趋近 $2a$。右图把割线与切线同时保留，是为了把“一个仍有宽度的测量”和“测量结果的极限”区分开。导数描述后者。



\paragraph{3. 零斜率不等于函数不变化}

把观察点移到 $a=0$。切线水平，而附近的函数仍向上弯曲。这里位置增量是 $h^2$，比一阶尺度 $h$ 更小；所以局部一阶读数为零。判断一个量是否不变，还必须考察邻域、较高阶变化或额外条件，不能只看单点导数。



\subsection*{定义、命题与证明}

\paragraph{定义：某一点的导数}

若极限存在，定义 $f'(a)=\lim_{h\to0,\ h\ne0}\dfrac{f(a+h)-f(a)}h$。这里先在非零步长处计算，再讨论读数的极限。

\paragraph{命题：平方函数的差商与误差}

对 $f(x)=x^2$ 及 $h\ne0$，差商为 $2a+h$；因此 $f'(a)=2a$，差商的绝对误差恰为 $|h|$。

\begin{proof}

展开 $(a+h)^2-a^2=2ah+h^2$，因 $h\ne0$，除以 $h$ 得 $2a+h$。

与候选极限 $2a$ 相减后仅剩 $h$。对任意正容差，只要步长的绝对值更小，读数误差也更小；这给出了极限所需的控制，而非在 $h=0$ 处进行除法。

\end{proof}
